\documentclass[conference]{IEEEtran}

\usepackage{amsmath,amssymb,amsfonts,amsthm}
\usepackage{booktabs}
\usepackage{tabularx}
\usepackage{array}
\usepackage[hidelinks]{hyperref}
\usepackage{url}

\hypersetup{
  pdftitle={WayBill: Receiver-Sealed Origin Labels and Odometer-Backed Settlement for Disclosure-Minimizing Road-Usage Charging},
  pdfauthor={Anonymous for blind review},
  pdfsubject={Privacy-preserving road-usage charging},
  pdfkeywords={road-usage charging, location privacy, zero-knowledge proofs, trusted sensing}
}

\theoremstyle{plain}
\newtheorem{theorem}{Theorem}
\theoremstyle{definition}
\newtheorem{definition}{Definition}

\newcommand{\system}{\textsc{WayBill}}
\newcommand{\ruc}{RUC}
\newcommand{\rmax}{r_{\max}}
\newcolumntype{Y}{>{\raggedright\arraybackslash}X}

\begin{document}

\title{\system: Receiver-Sealed Origin Labels and Odometer-Backed Settlement for Disclosure-Minimizing Road-Usage Charging}

\author{\IEEEauthorblockN{[Author names removed for blind review]}
\IEEEauthorblockA{[Affiliation removed for blind review]}}

\maketitle

\begin{abstract}
Disclosure-minimizing road-usage charging can avoid disclosing raw trajectories by computing
fees locally and proving correct payment. Correctness alone does not establish
completeness when a vehicle can label an interval as unavailable or withhold a
whole charging period. We present \system, a conditional settlement design
that assigns acquisition-label admissibility to a process-isolated receiver signer while
keeping tariff selection and bill computation outside that trusted process.
The signer seals acquisition-supplied fix records, monotone odometer readings,
and an admissibility bit derived from each origin-fix status. The vehicle proves in zero
knowledge that every reported interval follows a committed tariff or an
odometer-backed maximum-rate fallback. At month close, signed odometer distance
not covered by accepted non-overlapping periods is reconciled at the same
maximum rate. We prove that fallback and period withholding cannot reduce the
bill relative to acquisition-supplied, receiver-sealed tariff labels, conditional on correct
receiver attestations, proof soundness, and monthly odometry. The guarantee
does not cover receiver compromise, odometer tampering, or GNSS relay. A
receipted 25-fix Groth16 run of the V6 prototype contains 246,255 constraints
and 22 public signals. It produced an 803-byte JSON proof encoding in 3.65 s and verified in
217 ms with exact circuit and reference agreement. Endpoint receipts further
record replay rejection, public-signal tamper rejection, and month-close
reconciliation. These measurements validate one canonical execution path. The
pre-registered four-dataset trace evaluation of billing accuracy, scaling,
bounded displacement, metadata linkage, and handler behavior remains pending, and no
full-evaluation claim is made from the canonical receipt.
\end{abstract}

\begin{IEEEkeywords}
road-usage charging, location privacy, zero-knowledge proofs, trusted sensing,
accountable settlement, GNSS
\end{IEEEkeywords}

\section{Introduction}

Road-usage charging (RUC) can price distance, place, and time more directly
than a fuel tax. A direct Global Navigation Satellite System (GNSS) design,
however, gives the operator a detailed movement history. Private tolling
protocols instead compute charges locally and use commitments, proofs, or
audits to establish correct payment
\cite{Popa2009VPriv,Balasch2010PrETP,Meiklejohn2011Milo,Fetzer2020P4TC}.
This line of work has made fee correctness increasingly precise, yet a distinct
systems problem remains. A correct proof covers the records it contains. It
does not by itself show that every chargeable interval or period was submitted.

Completeness is commonly supported by infrastructure observations such as
gates, cameras, or mobile checks. A recent survey organizes these privacy and
enforcement tradeoffs, while a deployment case study shows that protocol-level
privacy alone need not remove linkability from the surrounding collection
infrastructure \cite{Jolfaei2023Survey,Jolfaei2024Brisbane}. The problem is
therefore broader than hiding coordinates inside a proof. A settlement design
must bind who decides whether location data are usable, account for periods
that never reach the verifier, and disclose the remaining account-level
metadata explicitly.

\system\ studies this problem for vehicles that already contain an enrolled
sensing unit and a calibrated monotone odometer. Its central design choice is
to separate origin-label admissibility from billing. A process-isolated receiver signer
seals fixes, odometer readings, and origin-label admissibility decisions. The normal
vehicle process proves tariff membership and bill arithmetic, while the
charger checks non-overlap and reconciles unreported monthly distance. This
boundary remains conditional. Navigation-message authentication can
authenticate transmitted navigation data, but it does not prove that a
receiver occupied a physical location or defeat every relay attack. Galileo
Open Service Navigation Message Authentication (OSNMA) status is consequently
an acquisition input to \system, not a complete location oracle
\cite{GalileoOSNMA2025,Psiaki2016GNSS,Sastry2003Location}.

The resulting argument is narrow but operationally useful. Receiver-sealed
validity prevents a normal prover from choosing a cheaper outage path, and
odometer-backed month close removes the discount from withholding a period
under stated sensing assumptions. The contributions are as follows.

\begin{enumerate}
  \item We define a validity-bound settlement interface in which a
  process-isolated receiver signer derives and commits origin-label
  admissibility with each fix. The normal
  prover cannot replace fallback bits, roots, tariffs, policy profiles, or
  verification keys after acquisition.
  \item We combine interval fallback with odometer-backed month accounting and
  prove a conditional no-discount result. Relative to
  acquisition-supplied, receiver-sealed tariff labels, neither declaring fallback nor
  withholding complete periods can lower the charge when the receiver,
  odometer, and proof assumptions hold.
  \item We implement the end-to-end V6 path and pre-register a five-stage
  evaluation over four mobility datasets. Immutable receipts bind code, data,
  container, policy, and artifacts. This draft reports only the canonical
  differential receipt and reserves all full-dataset claims for the strict
  server aggregate.
\end{enumerate}

These contributions separate relation correctness, conditional completeness,
and disclosure minimization. We make no universal deployment claim, no claim
of physical location truth, and no information-theoretic route-privacy claim.

\section{Related Work}

\subsection{Private tolling and completeness}

VPriv computes location-based fees while limiting disclosure through local
computation, commitments, and audits \cite{Popa2009VPriv}. PrETP embeds local
fee computation in an optimistic payment protocol and uses external evidence
to challenge committed positions \cite{Balasch2010PrETP}. The Phantom
Tollbooth addresses collusion and conceals spot-check locations during audits
\cite{Meiklejohn2011Milo}. P4TC develops provable privacy and security
properties for practical DSRC toll collection \cite{Fetzer2020P4TC}. The
broader literature spans centralized, distributed, and hybrid designs with
different observation and trust assumptions \cite{Jolfaei2023Survey}.
PriPAYD similarly demonstrates the relevance of local computation and trusted
on-board usage data outside toll collection \cite{Troncoso2007PriPAYD}.
\system\ does not subsume these protocols. It asks whether an enrolled receiver
and monthly odometry can replace part of the recurring observation function in
a GNSS setting, at the cost of a stronger trusted-sensing boundary.

The closest abstraction is a maximum-rate deposit followed by verified
discounts, which already gives the theorem's inequality. Our narrower systems
contribution is to make discount conditions auditable through receiver sealing,
proof binding, non-overlap, and an append-only month ledger. Attested on-board
meters or random roadside checks offer different assurance/availability
tradeoffs; we do not claim process isolation dominates them.

\subsection{Protocol privacy and deployment leakage}

Cryptographic privacy can be weakened by identifiers, account workflows,
receipts, and exception handling outside the core protocol. The Brisbane case
study demonstrates why the surrounding collection infrastructure belongs in
the privacy analysis \cite{Jolfaei2024Brisbane}. \system\ follows this lesson
by declaring the complete public settlement tuple and evaluating linkage under
several release views. It does not equate hidden coordinates with an anonymous
transaction.

\subsection{Authenticated navigation and zero-knowledge location}

Galileo OSNMA authenticates navigation messages received through the Galileo
Open Service \cite{GalileoOSNMA2025}. It does not establish the physical
presence of a receiver, so \system\ treats relay and meaconing as residual
sensing risks. Location-verification and GNSS-spoofing research make the
separate physical-presence problem explicit
\cite{Sastry2003Location,Psiaki2016GNSS}. Recent zero-knowledge proof-of-location work studies
territorial compliance for vehicle subsidies and taxation
\cite{Bogdanov2025ZKLocation}. Our relation instead binds a sequence of
receiver-sealed acquisition records to tariff membership, integer billing, a
period ledger, and month-close settlement. These systems solve related but
different tasks, and we make no superiority claim.

\section{System and Threat Model}

\subsection{Parties and trust boundary}

The \emph{tariff authority} publishes a versioned cell-to-zone table, zone
rates, a maximum rate $\rmax$, and a policy profile. The \emph{receiver signer}
is a process isolated from the normal prover. Here ``independent'' denotes
process and credential separation, not independent hardware, operator, or
sensor verification. An authorized acquisition path
appends fixes with a strict status enum. The signer, rather than the caller,
derives each outgoing interval bit under the profile-pinned
\texttt{origin-osnma-authenticated-v1} rule. The
\emph{prover} reads sealed period material and constructs a proof. The
\emph{charger} enrolls device public keys, pins policy profiles and verification
keys, accepts period statements, and closes monthly ledgers.

For each fix $f_i$, the receiver log contains a device and period identifier,
sequence number, authenticated GNSS time, cell label, odometer reading, nonce,
and device signature. For every nonterminal fix it also records
$v_i\in\{0,1\}$, indicating whether the profile admits the origin cell label for
pricing the next interval. This policy bit is not a physical-position or
continued-presence assertion. The normal prover has no API field with which to replace
$v_i$ or submit its own root. The receiver signer rejects sequence, time,
odometer, and log-epoch rollback. It also rejects duplicate period sealing and a log whose
bytes no longer match its stored hash. Before a new seal becomes durable, the
signer compares its validated local tail with an authenticated charger head and
atomically extends that per-device chain by one epoch. A stale snapshot, deleted
tail, epoch gap, or conflicting same-epoch attestation is rejected.

The trusted base includes the enrolled key and acquisition logic. That logic
supplies GNSS time, cell label, status, odometer reading, nonce, and monthly
odometer boundaries; the signer validates syntax, ordering, profile rules, and
chain continuity but does not parse OSNMA messages or independently recompute a
cell from a navigation solution. It also assumes odometer authenticity,
monotonicity, calibration accuracy, and continuity across maintenance. These
facts are software-signed, not remote or hardware attestations. The trusted
base excludes tariff selection and bill arithmetic. We do not address key
extraction, a maliciously enrolled receiver, undetectable physical odometer
tampering, ECU/CAN compromise, systematic under-recording, or a compromised
acquisition process; remote automotive attack surfaces motivate this exclusion
\cite{Checkoway2011Automotive}.

\subsection{Adversary and goals}

An adversarial prover may alter private witness values, substitute a tariff or
verification key, replay an old proof, rebind a statement to another period or
month, omit intervals or entire periods, or modify public signals. It may not
forge the receiver signature or the proof system. Signal relay and meaconing
are evaluated as residual sensing attacks, not ruled out cryptographically.

An accepted period should have the fee, distance, fallback count, roots,
profile commitment, and canonical time windows proved together. A withheld
period should not create an economic discount under the conditional month-close
model. The privacy goal is narrower. Raw fixes and interval charges are absent
from the proof-only request. The charger still learns the device-month link,
fee, distance, number of fallback intervals, time bounds, and commitments.

\section{WayBill Design}

\subsection{Design overview}

\system\ receives a sealed sequence of fixes for one charging period and
produces two linked outputs. The first is a zero-knowledge statement containing
the period totals and commitments required for verification. The second is an
immutable receipt that binds the accepted statement to receiver and charger
state. Four mechanisms form the settlement chain. A validity-bound receiver
root fixes which intervals may use ordinary tariff labels. A period relation
proves tariff membership and arithmetic. A verifier transition rejects replay
and overlap. Monthly reconciliation charges odometer distance absent from all
accepted periods. Figure~\ref{fig:flow} shows how responsibility passes between
these mechanisms.

\begin{figure*}[t]
\centering
\fbox{\begin{minipage}{0.94\textwidth}
\textbf{Acquisition and settlement flow.}
(1) An authorized acquisition component sends unsigned fixes with strict
status labels to a device-scoped receiver signer.
(2) The signer validates monotonicity, derives $v_i$ from the origin-fix
status, and commits each $f_i$ together with $v_i$ or a false terminal
sentinel. It seals the append-only log, signs the root metadata, and obtains an
exact compare-and-swap confirmation from the charger before committing the
local row.
(3) The vehicle prover obtains the sealed fixes, admissibility vector, and root
attestation, then computes tariff membership, fees, and a Groth16 proof.
(4) The charger verifies the profile-pinned proof, exact public-signal vector,
receiver attestation, period uniqueness, and non-overlap.
(5) Month close bills attested odometer distance not covered by accepted
periods at $\rmax$.
\end{minipage}}
\caption{The receiver owns the status-derived admissibility decision and root. The normal prover
owns neither. Pricing remains outside the receiver process.}
\label{fig:flow}
\end{figure*}

\subsection{Binding origin-label admissibility}

Let $H$ denote the circuit's Poseidon chaining construction. For a nonterminal
fix, the receiver and circuit compute
\[
 c_i=H(\mathsf{DAC},\mathsf{profile},\mathsf{period},i,t_i,x_i,y_i,
       o_i,n_i,v_i).
\]
where $\mathsf{DAC}$ identifies the data-acquisition contract,
$\mathsf{profile}$ and $\mathsf{period}$ bind the policy and period,
$(t_i,x_i,y_i,o_i,n_i)$ are the signed time, cell, odometer, and nonce fields,
and $v_i$ is the profile-derived origin-label admissibility bit. Here
$(x_i,y_i)$ are discrete cell labels, not receiver-signed sub-cell physical
coordinates.
The final fix uses the same tuple with a fixed false sentinel. Chaining all
$c_i$ yields $R_{\mathrm{fix}}$. Consequently, changing a validity decision
changes the root used by both the proof and the receiver's Ed25519 software
attestation. This
removes an interface in which signatures could cover fixes while leaving
fallback flags under prover control.

The signer attests the metadata tuple
\begin{align*}
(&\mathsf{receiver},\mathsf{device},\mathsf{period},\mathsf{epoch},
\mathsf{count},R_{\mathrm{fix}},\\
 &\mathsf{DAC},\mathsf{semantics},\mathsf{rule},\mathsf{charger},
 \mathsf{loghash},\mathsf{prevhash}).
\end{align*}
The attestation does not sign a statement supplied by the prover. The charger
first verifies the V5 attestation signature and appends its hash only if the epoch and
previous hash extend the current device head. At settlement, it compares the
root, device, period, and device-attestation commitment to the proved public
statement and accepts only the exact previously anchored attestation. The
read-only preflight avoids spending proof-verification work on an unanchored
request. The write transaction repeats the check to preserve correctness under
concurrency.

\subsection{Proving the period relation}

For interval $i$, let $\Delta o_i=o_{i+1}-o_i$. The circuit enforces odometer
and time monotonicity, a maximum time gap, an odometer speed bound, Boolean
$v_i$, tariff Merkle membership, integer width limits, and the profile-pinned
maximum rate. Its fee is
\[
 q_i=\begin{cases}
 \Delta o_i\,r(z_i),&v_i=1,\\
 \Delta o_i\,\rmax,&v_i=0.
 \end{cases}
\]
where $\Delta o_i$ is the nonnegative odometer increment, $z_i$ is
the acquisition-supplied, receiver-sealed tariff cell, $r(z_i)$ is its committed rate, and
$\rmax$ is the maximum rate fixed by the policy profile. A signer-sealed tariff
cell is therefore an acquisition-supplied label, not independent
physical-position verification.
Thus a stationary vehicle pays zero even during an outage. The fallback uses
the same odometer delta as normal settlement. All products and totals are
range-checked under the frozen uint32/uint64/uint96 contract.

The public vector contains 22 fields that bind the fix, tariff, and
interval roots, the policy-profile commitment, the period and month identifiers
and bounds, the tariff version and device commitment, the totals, the fallback
count, and the physical and policy bounds. The proof-only endpoint uses the
profile commitment to resolve one authority-approved circuit, tariff root, and
verification-key hash. Extra raw fix fields are rejected by the request
schema.

\subsection{Reconciling a month}

Accepted windows are non-overlapping for each device. At month close, the
receiver signs boundary readings $(o_{M,0},o_{M,1})$. If accepted periods cover
distance $D_A$, the charger computes
\[
 D_U=\max(0,o_{M,1}-o_{M,0}-D_A),\qquad Q_U=D_U\rmax.
\]
where $o_{M,0}$ and $o_{M,1}$ are the signed month-boundary odometer
readings, $D_A$ is the distance in accepted non-overlapping periods, $D_U$ is
the remaining distance, and $Q_U$ is its reconciliation charge.
An inconsistent excess $D_A$ is rejected outside a named tolerance. A missing
month attestation creates a fail-closed administrative exception rather than a
distance estimate. Collection policy for that exception is outside the
cryptographic theorem; it must not be treated automatically as fraud. Before
local persistence, the signer submits the exact
signed boundary to the authenticated charger. The first enrolled boundary is a
bootstrap trust point. Thereafter the charger accepts only the next calendar
month with $o_{M,0}$ equal to the preceding $o_{M-1,1}$. Exact retries are
idempotent, while different bytes for an anchored device-month are rejected.
This closes stale-local-database rollback at the cost of making charger
availability a prerequisite for sealing a month.

\subsection{Verifier state transition}

For each statement, the charger resolves the pinned profile, tariff root,
circuit, and verification key. It rejects an invalid or unanchored receiver
attestation, any mismatch with the exact public vector, a repeated
device-period key, or an overlapping window. Only after all checks pass does a
transaction record the period and its immutable receipt. When the month is
ready to close, the charger validates the boundary attestation against the
anchored month chain, computes $D_U$ and $Q_U$, and records a close receipt. An
invalid month boundary creates an administrative exception. No failed check
produces a billable period or advances the month ledger.

The prototype does not implement adjudication, but a deployable ledger needs an
append-only dispute state machine:
$\mathsf{accepted}\rightarrow\mathsf{notified}\rightarrow\mathsf{contested}
\rightarrow\mathsf{evidenceOpened}\rightarrow\mathsf{adjudicated}
\rightarrow\{\mathsf{adjusted},\mathsf{closed}\}$. A timely contest should
suspend penalties on the disputed amount. Evidence opens in stages---accounting
tuple, then endpoints, and raw trajectory only when necessary---under an
independent appeals authority. A correction appends a signed adjustment or
refund receipt and never overwrites the original. Deadlines, burden of proof,
refund interest, and access restrictions are policy obligations rather than
claims of the current implementation.

\section{Scoped Security Analysis}

\begin{definition}[Sealed-label reference bill]
For the subset of month distance being compared, associate each odometer
distance $d_i$ with a counterfactual acquisition-accepted tariff label $z_i$ and rate
$0\le r(z_i)\le\rmax$. The reference bill is
$Q_{\mathrm{auth}}=\sum_i d_i r(z_i)$. This definition concerns labels accepted
by the receiver acquisition boundary. Here $d_i$ is an odometer distance,
$z_i$ is its acquisition-accepted tariff label, and $r(z_i)$ is the committed rate.
The definition is not a physical-location oracle. A genuinely unavailable
interval for which no such reference label exists is outside this comparison
and is charged according to the fallback policy without a claim about its
unobserved zone rate.
\end{definition}

\begin{theorem}[Successful-close conditional economic conservatism]
\label{thm:withholding}
Assume (i) receiver root and month-boundary attestations are unforgeable,
(ii) accepted V6 proofs are sound and profile-pinned, (iii) monthly odometer
distance is correct and covers the considered month, and (iv) accepted periods
are non-overlapping and their accounted distance is subtracted exactly once.
Then replacing any reference-labeled interval by fallback, or withholding
any set of period statements and closing the month, cannot reduce the bill
below $Q_{\mathrm{auth}}$. The inequality is strict for positive withheld
distance whose sealed-label rate is below $\rmax$.
\end{theorem}

\begin{proof}
For an accepted interval with $v_i=1$, the circuit charges
$d_i r(z_i)$. With $v_i=0$, it charges $d_i\rmax\ge d_i r(z_i)$. Every distance
not represented by a non-overlapping accepted period remains in $D_U$ and is
charged at $\rmax$, again no less than its reference rate. Summing these
disjoint contributions proves the inequality. Strictness follows when
$d_i>0$ and $r(z_i)<\rmax$ for at least one replaced contribution.
\end{proof}

Equivalently, after a successful month close the mechanism starts from charging
all covered month distance at $\rmax$ and subtracts only the verified discounts
$d_i(\rmax-r(z_i))$ for accepted ordinary-rate intervals. The algebraic
monotonicity is immediate; the systems contribution is binding those discounts
to sealed labels, non-overlapping periods, and append-only state. The theorem
makes no claim for an indefinitely missing month attestation, which remains an
administrative exception requiring the dispute process above.

Theorem~\ref{thm:withholding} is intentionally not called ``physical tariff
truth.'' If a compromised receiver signs a false cell, the proof faithfully
bills that false sealed label. If odometry understates distance, month
close inherits the error. OSNMA authenticates navigation messages, whereas the
receiver's actual presence and resistance to delayed relay require additional
sensing and operational controls. These are assumption failures, not attacks
solved by the settlement circuit.

\paragraph{Proof and replay binding}
The charger compares the full public vector against the submitted statement,
uses only the profile-pinned verification key, checks the validity-bound root
attestation, rejects duplicate $(\mathsf{device},\mathsf{period})$ keys, and
rejects overlapping device windows. A valid proof detached from its original
statement therefore fails signal equality or profile resolution.

\paragraph{Privacy boundary}
We do not provide a simulation proof for an application-level privacy notion.
Instead, we declare the charger view and measure how uniquely it identifies a
period in a candidate pool. Hash commitments hide neither low-entropy public
totals nor linkability induced by the billing account. Bucketed release is a
policy mitigation with an explicit accuracy/leakage tradeoff.

\section{Implementation}

\subsection{Prototype and receipt discipline}

The V6 reference is implemented in Python and Circom. Its commitment and Merkle
constructions use the circuit-oriented Poseidon hash \cite{Grassi2021Poseidon}.
Groth16 is used for its
small proofs and mature tooling, not as a novel proof-system contribution
\cite{Groth2016}. The canonical policy is \texttt{ruc-demo-v9}. It pins tariff v9, the
validity-bound 25-fix circuit, maximum rate, integer semantics, and verification
key. The checked-in direct setup is explicitly test-only. Production requires
an independently governed ceremony or a different setup strategy.

The receiver signer is a device-scoped FastAPI process backed by owner-only
SQLite state. Its acquisition route requires a write bearer token, accepts only
unsigned measurement fields, derives rather than accepts the origin-label admissibility
bits, creates the receiver signatures itself, and seals each period and monthly
odometer boundary once. A separate read bearer token protects prover access to
sealed trajectories. Root-attestation V5 pins the validity rule and charger
domain and includes the previous attestation hash. The signer uses an
authenticated TLS channel to extend a Charger-side compare-and-swap head before
persisting the local row. An exact one-link-ahead retry recovers a lost response
or local commit. Monthly-attestation V2 also binds the intended charger domain
and is online-anchored under exact-retry and conflict-rejection semantics.
The signer pins the canonical profile itself, enforces its validity window and
fix cap, and opens its key, database, and profile with fail-closed ownership and
mode checks. Its SQLite schema is versioned and rejects injected triggers,
views, unknown objects, and unsupported legacy layouts. The charger stores
enrollments, accepted periods, month state, immutable receipts, and audit
events transactionally. The formal server configuration uses PostgreSQL,
administrator and device bearer credentials, and one application worker.
SQLite remains a local-test backend.
Formal stage-result V2 records include the job hash and recursive artifact
manifest. Attempt and final gate receipts are hashed. Merge revalidates files
from disk rather than trusting a copied \texttt{passed} value.

\section{Evaluation Methodology and Current Evidence}

\subsection{Canonical-path validation}

Table~\ref{tab:v6} reports the checked-in receipt generated on an arm64
development host with snarkjs. Timing comes from one artifact-validation run
and is not a population estimate. The receipt explicitly records the test-only
setup.

\begin{table}[t]
\centering
\caption{Canonical validity-bound V6 artifact receipt.}
\label{tab:v6}
\small
\begin{tabular}{lr}
\toprule
Metric & Receipt value \\
\midrule
Fixes / intervals & 25 / 24 \\
Constraints & 246,255 \\
Public signals & 22 \\
Fallback intervals & 3 \\
Witness generation & 0.408 s \\
snarkjs proving & 3.646 s \\
snarkjs verification & 0.217 s \\
Proof JSON & 803 bytes \\
Circuit/reference vector & exact match \\
\bottomrule
\end{tabular}
\end{table}

The same proof passed an in-process endpoint-logic gate using FastAPI
TestClient and SQLite test mode; this is not deployed server evidence.
Re-submission was
rejected with HTTP 409. Changing the public fee signal was rejected with HTTP
400. The accepted period covered 2,400 m and 3,600 cents. A signed 3,000 m month
boundary caused exactly 600 m to be reconciled at five experimental cents/m,
for a 3,000-cent
supplement. The endpoint receipt also records that its root attestation came
from a separate receiver process, names the sealed-log SHA-256, and records
the exact V5 chain head anchored before proof acceptance. It separately binds
the monthly attestation SHA-256, confirms exact-retry idempotence, and records
rejection of a conflicting validly signed boundary for the same device-month.

\subsection{Research questions and pre-registered protocol}

The full evaluation is organized around five research questions. These trace
datasets contain neither real OSNMA verification states nor independently
measured vehicle odometers; the experiment is mechanism sensitivity, not a
certification of receiver correctness, anti-tamper odometry, or physical
location truth. RQ1 asks whether V6 preserves trace-derived reference billing
accuracy under synthetic independent and bursty unavailability, synthetic
odometer error, and whole-period withholding. RQ2 measures
how circuit depth and period size affect compilation, setup, proving,
verification, memory, and JSON proof-encoding size. RQ3 quantifies the billing
effect of bounded cell/trajectory displacement across tariff granularities; it
is not a signal-level relay or meaconing experiment. RQ4 measures
identity linkage under exact, fallback-augmented, bucketed, and selectively
opened public views. RQ5 tests end-to-end CLI proof verification and the
in-process FastAPI handler/SQLite ledger under valid, tampered, and replayed
concurrent workloads. It does not measure HTTPS/Uvicorn/PostgreSQL deployment
capacity.

The protocol uses T-Drive, GeoLife, Porto, and Rome mobility traces. It applies
five fixed seeds and does not cap trips, raw points, inputs, or fixes. S1
compares a trace-derived position reference, the earlier V5 relation, V6, and V6
without month close. The earlier private tolling systems provide design
baselines for local computation, commitment auditing, and spot-check
enforcement \cite{Popa2009VPriv,Balasch2010PrETP,Meiklejohn2011Milo}. They are
not relabeled as executable V6 baselines because their trust and observation
models differ.

Formal jobs will run in a frozen Linux x86-64 container with Python 3.12, Circom
2.1.9, snarkjs 0.7.6, and the policy-pinned artifacts. Each setup or proving job
receives eight CPU cores and 32 GiB of memory. The largest handler stage
receives 32 CPU cores and 64 GiB. The target starts with at least 500 GiB of
free workspace and the exact 4.83 GB powers-of-tau file. These are scheduler
allocations and admission thresholds. Processor model and operating-system
details will be read from the target-host receipt after the server run.

\begin{table}[t]
\centering
\caption{Pre-registered evaluation stages and current evidence boundary.}
\label{tab:formal-status}
\small
\begin{tabularx}{\columnwidth}{@{}lYY@{}}
\toprule
Stage & Primary evidence & Draft status \\
\midrule
S1 & Billing accuracy and completeness & Server aggregate pending \\
S2 & Proof cost and scaling & Canonical V6 receipt only \\
S3 & Bounded displacement sensitivity & Server aggregate pending \\
S4 & Metadata linkage & Server aggregate pending \\
S5 & CLI verifier and in-process handler & Server aggregate pending \\
\bottomrule
\end{tabularx}
\end{table}

S1 will report vehicle or user cluster-bootstrap 95\% confidence intervals with
5,000 replicates and treats every rejected period as fail closed. Its
month-close ablation removes the complete close mechanism. S2 will record warm-up
trials separately and report the median, 95th percentile, and 99th percentile
for measured trials. With ten trials, P95 and P99 are descriptive nearest-rank
maxima rather than stable tail estimates. S3 freezes four tariff bucket sizes
and three bounded-displacement radii. S4 separates four disclosure views. These are exact totals, exact
totals with fallback metadata, bucket-only totals, and bucket-only totals with
opened cells. Bucket views contain no duplicate exact totals. Gallery and query
training and test identities are disjoint. Gallery and query records use the
same held-out identities but disjoint periods, and a one-period temporal gap is
enforced. S5 will use distinct real Groth16 bundles rather than repeatedly
verifying one cached proof.

These choices avoid two evaluation shortcuts. A bootstrap over rows would
overstate precision when periods from one driver are correlated, so the
resampling unit is the original vehicle or user. Retaining exact totals beside
bucketed values would make a bucketed treatment observationally identical to
the exact view, so each view has its own strict schema.

\subsection{Evidence boundary}
Table~\ref{tab:v6} is the only quantitative proof-performance result claimed
before the formal run. The four-dataset S1--S5 results will be inserted only
from the strict aggregate after every expected job has a successful immutable
attempt. A timeout, out-of-memory event, rejected period, or missing dataset is
retained in hash-bound \texttt{failure-summary.json} and
\texttt{scientific-aggregate.json}. The latter binds selected attempts,
parameters, stage metrics, per-seed rows, and median/min/max cross-seed summaries.
No result is replaced by a smaller experiment or copied manually from a log.

% RESULT_SLOT_S1: billing accuracy and completeness aggregate
% RESULT_SLOT_S2: proof cost and scaling aggregate
% RESULT_SLOT_S3: bounded-displacement and tariff-sensitivity aggregate
% RESULT_SLOT_S4: metadata-linkage aggregate
% RESULT_SLOT_S5: CLI-verifier and in-process-handler aggregate

\section{Discussion}

The contribution of \system\ is a settlement boundary rather than a new claim
that cryptography can establish physical location. Earlier private tolling
systems combine local fee computation with external observations or
protocol-specific audits. \system\ instead assigns three kinds of evidence to
three mechanisms. The receiver seals acquisition records and profile-defined
label admissibility, the proof system
attests pricing arithmetic, and the monthly odometer attests total distance.
This decomposition makes withholding economically conservative under the
theorem's assumptions while preserving a proof-only, but account-linkable,
interface for ordinary settlement. It relocates tariff selection and bill
arithmetic outside the acquisition boundary, but does not by itself reduce the
acquisition component's sensing privilege. Whether this trade is preferable to roadside
observation depends on deployment governance, hardware assurance, and the
acceptable metadata view.

\subsection{Residual sensing and metadata risk}

Earlier trajectory sweeps are retained as versioned, non-V6 empirical context,
not merged into the canonical cryptographic receipt. They show why the theorem
is scoped. Proof-consistent relay can preserve a nonzero billing residual, and
exact fee--distance--time tuples can have an anonymity set of one in small pools.
Bucketed totals improve this empirical anonymity measure, while opening endpoint
fixes during a dispute can sharply reduce it. Consequently, deployment should
combine signal-diversity controls, conservative fallback, bucketed public
release where policy permits, and accounting-level dispute openings before raw
endpoint disclosure.

\subsection{Limitations}

The independent signer is a process boundary, not a remotely attested hardware
implementation. Separate acquisition and prover tokens reduce privilege, but a
compromised acquisition component can still falsify a status label or fix. The
derivation rule does not turn software input into hardware truth. Token and key
management require rotation, hardware protection, and audit. Online Charger
anchoring detects local database rollback, but it does not protect a stolen
receiver private key or a host already compromised under the signer UID. A
production deployment should move signing and monotonic state into attested
hardware. The PostgreSQL charger path has transaction/restart/concurrency tests,
but the paper does not claim multi-host failover or multi-worker throughput
before target-host stress evidence exists. The canonical proof uses a test-only
Groth16 setup. The 25-fix profile is an artifact-scale configuration, not
evidence of month-scale client throughput. The current theorem assumes a
correct odometer and a completed month attestation. Administrative handling of
missing month close remains a governance problem. Finally, the public
settlement tuple is linkable and potentially identifying, and we provide an
empirical leakage analysis rather than a composable privacy proof.

\subsection{Operational and ethical implications}

\paragraph{Billing errors and recourse}
A private bill is not automatically a fair bill. Receiver faults, tariff-map
errors, or conservative fallback can overcharge an honest driver. A deployment
therefore needs notice, bounded retention, an appeal deadline, and staged
dispute openings. It should disclose accounting tuples first and raw endpoints
only when necessary. Administrative handling of a missing month attestation must be
defined by public policy and cannot be silently interpreted as evidence of
fraud. Neither the prototype nor Theorem~\ref{thm:withholding} authorizes an
automatic penalty.

M4 cannot establish distributional fairness because its traces lack reliable
social and device-quality strata. Deployment work should report false-invalid
overcharges, affected-account share, longest fallback run, and correction delay
across difficult environments; disputed maximum-rate amounts should remain
provisional and receive signed refunds after a confirmed fault.

Governance also lies outside the prototype: tariff authorities own rates/maps,
vendors own signer updates/keys, operators own availability/retention, drivers
receive notice and contest, and a separate authority audits and adjudicates.
Process separation alone does not create organizational independence.

\paragraph{Availability}
The rollback defense deliberately makes the Charger an online dependency for
new seals. A network, TLS, or database outage pauses sealing rather than falling
back to an unanchored local signature. Acquisition may retry the exact period
after recovery. A timeout after remote commit is recoverable only when the same
bytes reproduce the one-link-ahead head. This fail-closed choice prevents silent
forking but creates an operational availability and denial-of-service risk.

\paragraph{Data governance}
Mobility traces are sensitive even when pseudonymized. The four-dataset formal
protocol records provenance and hashes, but restricted raw records are not
distributed with the repository. Published aggregates must satisfy provider
terms and exclude row-level trajectories, device keys, receiver logs, and
selective-disclosure payloads. An operational charger should minimize retention
of linkable public statements after billing and appeal obligations expire.

\paragraph{Artifact availability}
The repository includes the V6 circuit and wrapper, Python reference,
receiver-signer and charger services, policy/tariff artifacts, verification
key, tests, formal orchestrator, V2 receipts, and the small canonical proof
bundle. Large R1CS/WASM/zkey files, the test PTAU, and restricted datasets are
referenced by immutable hashes and distributed or mounted separately. The
direct Groth16 setup is marked test-only and is insufficient evidence of a
production ceremony. The artifact guide gives focused and full verification
commands and distinguishes the cryptographic receipt from unavailable formal
dataset stages.

\section{Conclusion}

\system\ demonstrates a precise separation between receiver-sealed
origin-label admissibility and publicly verified pricing. Binding that bit into the sealed root
removes the prover's ability to self-declare an outage, while odometer-backed
fallback and month close give a conditional no-discount result for
acquisition-supplied, receiver-sealed labels. The result can reduce raw trajectory disclosure
and some roadside dependence, but only under explicit sensing and odometry
assumptions. Version-pinned artifacts, process-isolated-signer receipts, and
fail-closed experiment manifests make those boundaries auditable.
The canonical receipt validates exact circuit and reference agreement, replay
rejection, public-signal tamper rejection, and month-close reconciliation. It
does not replace the four-dataset formal evaluation. \textbf{Future Work.} We
will complete the preregistered S1--S5 server run, report target-host resource
and failure distributions, evaluate attested receiver hardware, and study
release mechanisms that reduce settlement linkability without weakening
auditability.

\bibliographystyle{IEEEtran}
\bibliography{references}

\end{document}
